\documentclass[UTF8,a4paper,11pt,fontset=fandol]{ctexart}
\usepackage[left=17mm,right=17mm,top=17mm,bottom=18mm]{geometry}
\usepackage[table]{xcolor}
\usepackage{tikz}
\usetikzlibrary{arrows.meta,positioning,calc}
\usepackage{tcolorbox}
\usepackage{array,tabularx,booktabs}
\usepackage{fancyhdr}
\usepackage[hidelinks]{hyperref}

\definecolor{Ink}{HTML}{14253B}
\definecolor{Muted}{HTML}{455D73}
\definecolor{Rule}{HTML}{CCD7E2}
\definecolor{Blue}{HTML}{1D568D}
\definecolor{Mint}{HTML}{176B5E}
\definecolor{Amber}{HTML}{8A4A15}
\definecolor{Red}{HTML}{9A3B3E}
\definecolor{Ice}{HTML}{EFF5FB}
\definecolor{Leaf}{HTML}{EAF6F2}
\definecolor{Sand}{HTML}{FBF0E1}
\definecolor{Rose}{HTML}{FAEDEE}
\renewcommand{\familydefault}{\sfdefault}
\hypersetup{pdftitle={AI NPC 框架},pdfauthor={Delivery 项目}}
\setlength{\parindent}{0pt}
\setlength{\parskip}{5pt}
\setlength{\tabcolsep}{8pt}
\renewcommand{\arraystretch}{1.32}
\setlength{\headheight}{15pt}
\raggedbottom
\pagestyle{fancy}
\fancyhf{}
\renewcommand{\headrulewidth}{0pt}
\renewcommand{\footrulewidth}{0.35pt}
\fancyfoot[L]{\small\color{Muted}Delivery \hspace{.45em}/\hspace{.45em} AI NPC 框架}
\fancyfoot[R]{\small\color{Muted}\thepage\ / 2}

\newcommand{\hero}[2]{%
  {\small\bfseries\color{Blue}DELIVERY\enspace /\enspace UE 5.8\enspace /\enspace 项目面试材料}\par
  \vspace{3mm}{\fontsize{31}{38}\selectfont\bfseries #1}\par
  \vspace{1mm}{\large\color{Muted}#2}\par\vspace{3mm}}
\newcommand{\sectionline}[1]{\par\vspace{2mm}{\large\bfseries #1}\par\vspace{.5mm}}
\newcommand{\fig}[2]{\par{\centering #1\par}\vspace{-.5mm}%
  {\centering\small\color{Muted}#2\par}\vspace{1mm}}
\newcommand{\aside}[2]{%
  \begin{tcolorbox}[colback=Ice,colframe=Rule,boxrule=.5pt,arc=2mm,boxsep=0pt,top=2.4mm,bottom=2.4mm,left=3.2mm,right=3.2mm,before skip=4pt,after skip=4pt]
  \textbf{\color{Blue}#1}\par\small #2
  \end{tcolorbox}}
\newcommand{\refline}[1]{\vspace{2mm}\par{\footnotesize\color{Muted}\textbf{工程参考}\enspace #1\par}}
\newcommand{\codeword}[1]{\texttt{#1}}
\newcommand{\celltitle}[1]{\textbf{\color{Ink}#1}}
\newcolumntype{Y}{>{\raggedright\arraybackslash}X}
\tikzset{
  box/.style={draw=Rule,fill=Ice,rounded corners=2.2mm,align=center,
    font=\small,text=Ink,inner sep=2mm,text width=3.2cm,minimum height=13mm},
  engine/.style={box,fill=Leaf,draw=Mint!42!white},
  backend/.style={box,fill=Sand,draw=Amber!42!white},
  issue/.style={box,fill=Rose,draw=Red!38!white},
  head/.style={font=\bfseries\small,text=Ink},
  label/.style={font=\footnotesize,text=Muted,align=center},
  arrow/.style={-{Stealth[length=2.2mm,width=1.5mm]},line width=.75pt,draw=Blue},
  returnpath/.style={arrow,densely dashed,draw=Mint},
  line/.style={line width=.55pt,draw=Rule}
}
\begin{document}
\color{Ink}

\hero{AI NPC 框架}{一句话怎样变成看得见的动作}
Delivery 是一款送快递游戏。玩家可以和街坊聊天、起哄或劝架。每个 NPC 从自己的视角判断接下来做什么，游戏服务器负责判定这件事能否发生。

\sectionline{从玩家输入到双方看见结果}
\fig{\begin{tikzpicture}[x=1cm,y=1cm]
  \node[label,anchor=west] at (-1.8,.98) {玩家端};
  \node[label,anchor=west,text=Mint] at (2.53,.98) {主机 UE};
  \node[label,anchor=west,text=Amber] at (6.84,.98) {AI 后端};
  \node[label,anchor=west,text=Mint] at (11.15,.98) {主机 UE};
  \node[box] (input) at (0,0) {\textbf{交谈与感知}\\文字、麦克风、现场事件};
  \node[engine] (agent) at (4.31,0) {\textbf{NPC Agent}\\整理证据，列出合法动作};
  \node[backend] (model) at (8.62,0) {\textbf{模型决策}\\选一个动作，生成台词};
  \node[engine] (gate) at (12.93,0) {\textbf{服务器复核}\\检查请求、目标和状态};
  \draw[arrow] (input)--(agent);
  \draw[arrow] (agent)--(model);
  \draw[arrow] (model)--(gate);
  \node[engine] (execute) at (12.93,-2.25) {\textbf{行为树与技能}\\移动、出拳、命中判定};
  \node[box] (view) at (8.62,-2.25) {\textbf{主机与客人}\\动作、表情、字幕和声音};
  \node[engine] (facts) at (4.31,-2.25) {\textbf{结果回到 NPC}\\命中、挥空、取消或完成};
  \draw[arrow] (gate)--(execute);
  \draw[arrow] (execute)--(view);
  \draw[returnpath] (execute.south) -- ++(0,-.38) -| (facts.south);
  \draw[returnpath] (facts)--(agent);
\end{tikzpicture}}{实线是执行路径，虚线把实际动作结果带回下一轮决策。}

\sectionline{谁负责什么}
\begin{tabularx}{\linewidth}{@{}p{34mm}Y@{}}
\toprule
\celltitle{游戏内组件} & \celltitle{主要职责}\\
\midrule
\textbf{Agent 与 Gateway} & Agent 保存各 NPC 的证据和当前计划。Gateway 请求后端，按请求与世界标识关联回包。\\
\textbf{模型与语音后端} & 主对话模型理解玩家说的话。主机请求台词配音，语音按句分段发送。\\
\textbf{行为树与技能} & 行为树负责走到目标附近、等待和收招。技能执行出拳，服务器按实际击中结算。\\
\bottomrule
\end{tabularx}

\vspace{3mm}
\aside{为什么要列出合法动作？}{模型只需在本轮可执行的选项里做选择。比如目标被遮挡或走不到，游戏会先筛掉对应的追打选项。目标在等待期间离开，服务器还会再校验一次。}

\newpage
\hero{一次居民交战}{从起哄、动手到停手，动作与台词如何配合}
在双人局里，主机为两个 NPC 分别做决策。玩家在场看到的命中、受击和字幕来自同一局游戏状态。

\sectionline{服务器如何处理同一轮行动}
\fig{\begin{tikzpicture}[x=1cm,y=1cm]
  \node[box] (hear) at (0,0) {\textbf{听见玩家发言}\\保留谁说、对谁说};
  \node[backend] (choose) at (4.31,0) {\textbf{选择行动与台词}\\看向、打一拳、追打};
  \node[engine] (start) at (8.62,0) {\textbf{复核后开始}\\检查身份、距离和路径};
  \node[engine] (end) at (12.93,0) {\textbf{结算实际结果}\\命中、挥空、受击中止};
  \draw[arrow] (hear)--(choose);
  \draw[arrow] (choose)--(start);
  \draw[arrow] (start)--(end);
  \node[label,anchor=west] at (0,-1.05) {新的现场事件};
  \draw[returnpath] (end.south) -- ++(0,-.83) -| (choose.south);
  \node[label,fill=white,inner sep=1.2mm] at (8.55,-1.68) {受击或目标变化后重新判断};
\end{tikzpicture}}{同一场争执可以在继续追打、逃跑和停手之间变化。}

\sectionline{文字、麦克风与语音回放}
\fig{\begin{tikzpicture}[x=1cm,y=1cm]
  \node[box] (mic) at (0,0) {\textbf{玩家电脑}\\麦克风在本地转文字};
  \node[engine] (host) at (4.31,0) {\textbf{主机}\\校验发言并请求决策};
  \node[backend] (voice) at (8.62,0) {\textbf{后端}\\生成台词与分段音频};
  \node[box] (clients) at (12.93,0) {\textbf{听者客户端}\\显示字幕并按时序播放};
  \draw[arrow] (mic)--(host);
  \draw[arrow] (host)--(voice);
  \draw[arrow] (voice)--(clients);
  \node[label,anchor=west] at (-1.8,-1.05) {文字输入从发言校验进入同一条链路。};
  \node[label,anchor=west] at (6.79,-1.05) {受击或轮次取消会停掉相应播放。};
\end{tikzpicture}}{每个玩家可本地识别语音，服务器统一安排 NPC 的决定和发言。}

\aside{战斗小模型放在哪里？}{在启用的试玩配置里，已开始的追打或逃跑可以交给本地小模型继续判断。玩家发言、初次挑拨和需要回答的对话走主模型。小模型超时或输出无效时回到主模型。}

\refline{\path{Document/Architecture.md} 的居民框架、行为树与语音章节。\par
\path{Source/Delivery/AI/Decision/DeliveryNpcAgentComponent.cpp}、\path{Tools/DeliveryAI/combat_model.py}。}
\end{document}
